\section{Conclusion and future work}
\label{sec:conclusion}
This study asked whether the benchmark scores by which prompt-injection detectors are chosen predict their behavior inside an LLM agent. We evaluated fifteen released detectors and two task-aware judges on replayed tool outputs from two agent benchmarks and on BIPIA, and audited the published training data of two detectors. Detection rankings transferred poorly: the detector that ranked first on BIPIA, with \BIPGTpr detection at 1\% FPR, caught \ADPGTpr of AgentDojo injections, and the rankings on the two agent benchmarks agreed only partly. What a detector was trained on explained where it worked. The BIPIA leader had been trained on full BIPIA inputs, seeing an attack's text without its surrounding tool output did not help, and the detector that led both agent benchmarks shared no data with them but had been trained on agent-style inputs, keeping its lead with held-out thresholds (\LoHZTpr). False-positive behavior on tool outputs, by contrast, was stable across agent benchmarks and cheap to measure.

For practitioners, a detector should be selected and calibrated on the agent's own tool outputs, with false positives counted as blocked tasks, a threshold set on held-out traffic, and a check of what the detector was trained on; LLM-free replay with differential labels and published training sources make these checks possible.

Future work should extend the evaluation in three directions: to adaptive attackers who know the detector, to tool outputs from deployed agents rather than simulated environments, and to closed commercial detectors whose training data cannot be audited. Larger benign sets would also tighten the estimates at the low false-positive rates that deployments need.
